\documentclass[12pt]{article}
\usepackage[T1]{fontenc}
\usepackage[utf8]{inputenc}
\usepackage{lmodern}
\usepackage{textcomp}
\usepackage{enumitem}
\usepackage{geometry}
\geometry{margin=1in}
\begin{document}
\begin{center}
Answer Key - Religious Studies - K-12 Level Assessment 24 \
\small (Answers are ordered exactly as in the question paper.)
\end{center}

\section*{Multiple Choice}
\begin{enumerate}[label=\arabic*.]
\item \textbf{Answer:} A
\\ \textit{Options:} A) Palestine and Babylonia; B) Babylonia and Europe; C) Palestine and Spain; D) Spain and France
\\ \textit{Note:} Palestine and Babylonia became the two main centers for Jewish development after the Bar Kochba revolt due to the significant Jewish populations and the establishment of important religious and scholarly institutions in these regions.
\item \textbf{Answer:} C
\\ \textit{Options:} A) Genshin; B) Dogen; C) Honen; D) Shinran
\\ \textit{Note:} Honen is recognized as the founder of Pure Land Buddhism in Japan, where he emphasized devotion to Amida Buddha and the practice of reciting the nembutsu as the means to attain enlightenment and rebirth in the Pure Land.
\item \textbf{Answer:} B
\\ \textit{Options:} A) Loving and forgiving; B) Wrathful and merciful; C) Judging and vengeful; D) Transcendent and immanent
\\ \textit{Note:} The composition portrays Marduk as both wrathful and merciful by illustrating his capacity for fierce judgment against chaos while simultaneously offering protection and guidance to his followers, reflecting the dual nature of divine authority in ancient Egyptian beliefs.
\item \textbf{Answer:} C
\\ \textit{Options:} A) Indonesia; B) China; C) Korea; D) Singapore
\\ \textit{Note:} Korea is home to the largest network of Confucian shrines, known as Seowon, which were established as centers for Confucian learning and worship, reflecting the significant influence of Confucianism on Korean culture and education.
\item \textbf{Answer:} C
\\ \textit{Options:} A) Kama; B) Samnyasin; C) Ashramas; D) Arthas
\\ \textit{Note:} The term "Ashramas" refers to the four stages of life in Hindu philosophy, which are Brahmacharya (student), Grihastha (householder), Vanaprastha (hermit), and Sannyasa (renounced), as outlined in the dharma texts of the classical period.
\end{enumerate}

\section*{True / False}
\begin{enumerate}[label=\arabic*.]
\setcounter{enumi}{5}
\item \textbf{Answer:} True
\\ \textit{Options:} A) True; B) False
\\ \textit{Note:} Janam-sakhis are indeed a recognized genre of Sikh literature that specifically narrate the life and birth stories of Guru Nanak, highlighting his significance as the founder of Sikhism.
\item \textbf{Answer:} False
\\ \textit{Options:} A) True; B) False
\\ \textit{Note:} The festival that marked the beginning of the Babylonian year is called Akitu, not Wag and Thoth.
\item \textbf{Answer:} False
\\ \textit{Options:} A) True; B) False
\\ \textit{Note:} The term 'Arhats' refers to individuals in Buddhism who have achieved enlightenment and liberation from the cycle of rebirth, not to those who possess great wealth and material success.
\item \textbf{Answer:} True
\\ \textit{Options:} A) True; B) False
\\ \textit{Note:} Charism refers to a special grace or spiritual gift bestowed by the Holy Spirit that enables individuals to contribute to the well-being and service of their community.
\item \textbf{Answer:} True
\\ \textit{Options:} A) True; B) False
\\ \textit{Note:} The Hindu Vedas are considered the most sacred texts in Hinduism, and their recitation and interpretation are primarily the responsibility of Brahmins, who hold the priestly position within the religious hierarchy.
\end{enumerate}

\section*{Long Form Response}
\begin{enumerate}[label=\arabic*.]
\setcounter{enumi}{10}
\item \textbf{Answer:} The term "Guru-Panth" is significant in Sikh traditions as it embodies the concept of community within the Sikh faith. It refers to the collective body of Sikhs who follow the teachings of the Gurus and work together to uphold the values of equality, service, and justice. This idea emphasizes the importance of unity and collaboration among individuals, highlighting that spiritual growth is not just a personal journey but a shared experience within the community. The Guru-Panth illustrates how Sikhs are encouraged to support one another and participate actively in communal activities, reinforcing the belief that a strong, interconnected community is essential for both individual and collective spirituality.
\\ \textit{Note:} correct because it accurately defines "Guru-Panth" as the collective body of Sikhs that emphasizes community, shared values of equality and service, and the interconnectedness of individual spiritual growth within the Sikh faith.
\item \textbf{Answer:} Sri Lakshmi is a central figure in the celebration of Divali, symbolizing wealth, prosperity, and good fortune. Her attributes include beauty, grace, and abundance, making her a beloved goddess among Hindus. During Divali, rituals such as the lighting of oil lamps, the decoration of homes, and the offering of sweets and prayers are performed in her honor to invite her blessings into households. The festival not only emphasizes the importance of material wealth but also highlights the spiritual aspect of gratitude and community spirit. Overall, the celebration of Divali and the veneration of Sri Lakshmi reflect cultural values of hope, renewal, and the pursuit of prosperity in all aspects of life.
\\ \textit{Note:} correct because it accurately identifies Sri Lakshmi as a symbol of wealth and prosperity during Divali, details the rituals performed in her honor, and connects the festival's significance to both material and spiritual aspects of life.
\end{enumerate}

\vfill
\noindent\textit{End of Answer Key}
\end{document}